\chapter{Wstęp}
\label{chapt:Wstep}

Podczas czytania książek uwaga czytelnika przykuwana jest przez sam tekst, jednak na ostateczne doświadczenie lektury wpływa wiele czynników, które na pierwszy rzut oka mogą wydawać się niepowiązane. Jednym z takich elementów jest muzyka towarzysząca czytaniu, może ona nadać ton, pogłębić odczuwany smutek lub dać poczucie lepszej immersji w świat dzieła. Celem pracy jest stworzenie systemu wspomagającego przeżywanie lektury poprzez automatyczny dobór muzyki na podstawie analizy emocjonalnej tekstu z wykorzystaniem uczenia maszynowego.

\section{Cel pracy}
Opracowanie i realizacja aplikacji webowej przeznaczonej do lektury e-booków, która w sposób
zautomatyzowany personalizuje doświadczenia użytkownika poprzez dobór podkładu muzycznego
pasującego do nastroju aktualnie odtwarzanego fragmentu tekstu. System umożliwi analizę treści w
czasie rzeczywistym, wykorzystując modele uczenia maszynowego do klasyfikacji emocjonalnej tekstu
i dynamicznego sterowania odtwarzaczem audio. Zastosowanie do implementacji rozwiązania
nowoczesnych technologii (React.js, FastAPI/Python), architektury REST API oraz kontenerów
Docker, co zapewni skalowalność systemu i ułatwi jego wdrożenie.

\section{Aspekt inżynierski}
Projekt oraz implementacja uniwersalnej aplikacji do odczytu e-booków umożliwiającej automatyczny
dobór muzyki ambientowej w zależności od aktualnie prezentowanej treści. Realizacja pełnego stosu
technologicznego (full-stack) systemu: warstwy prezentacji (technologia React.js lub alternatywna) -
zapewnia interfejs czytnika obsługujący formaty EPUB/PDF; logiki biznesowej (technologia
FastAPI/Python) - integruje potok przetwarzania języka naturalnego (NLP), w którym model Sentence
Transformer dokonuje ekstrakcji cech tekstu (embedding), a wytrenowany klasyfikator
XGBoost/LightGBM przypisuje mu określony nastrój; warstwy danych - optymalizacja wydajności
aplikacji poprzez udostępnianie informacji o utworach w relacyjnej bazie SQLite, natomiast plików
audio za pośrednictwem wydajnego serwera statycznego Nginx. Podejmowanie decyzji o zmianie
ścieżki dźwiękowej na podstawie analizy statystycznej wyjścia modelu (thresholding) oraz dokonanej
oceny poziomu emocjonalnego aktualnie odtwarzanego fragmentu tekstu (eliminacja zbyt częstych
zmian muzyki). Wdrożenie systemu z użyciem technologii Docker oraz środowiska Docker Compose
(izolacja usług i ułatwienie zarządzania zależnościami). Weryfikacja poprawności działania aplikacji z
wykorzystaniem testów jednostkowych, integracyjnych i funkcjonalnych E2E. Analiza wyników testów
oraz wnioski. 

\section{Wykorzystanie AI}
W trakcie realizacji pracy wykorzystano modele językowe Claude Sonnet (Anthropic) oraz Gemini (Google) jako wsparcie w procesie tworzenia oprogramowania i redagowania dokumentacji. Narzędzia te były używane w szczególności do: wyszukiwania i syntezy informacji na temat wykorzystywanych technologii (React, FastAPI, biblioteki uczenia maszynowego); wsparcia w roli asystenta programisty (\textit{pair programming}) przy pisaniu fragmentów kodu, które następnie były dogłębnie analizowane, testowane i w razie potrzeby modyfikowane. Nauki nowych technologii i wcześniej nieznanych wzorców projektowych. Generowania wstępnych szkiców dokumentacji technicznej kodu oraz propozycji scenariuszy testowych, które zostały zweryfikowane i dostosowane do rzeczywistych wymagań systemu. Cały kod źródłowy został napisany samodzielnie lub zredagowany po dogłębnym zrozumieniu jego działania.
